
This work measured the effect of fastText quality filtering on demographic-occupation association structure in pretraining corpora derived from DCLM-POOL at quick-run scale (~50,000 documents per configuration). Two corpus-level findings were established with statistical significance: (1) fastText filtering reduces conditional entropy H(occupation|demographic) by 22.41\% between the 10th and 90th percentile thresholds (Spearman rho = -1.0, p = 1.4 x $10^-24$; bootstrap 95\% CI [-1.154, -0.330]), and (2) conditional log-odds of demographic-occupation co-occurrences increase monotonically with filtering intensity across 1,800 pairs (Spearman rho = 1.0, p = 1.4 x $10^-24$; mean log-odds from 0.697 to 2.976). A proxy-scale model training pilot did not yield statistically significant evidence of corpus-to-model signal propagation (rho = 0.357, p = 0.432), though a negative control showed detectable differences. The downstream fairness benchmark evaluation was not executed.

These results characterize fastText filtering as a mechanism that restructures demographic associations at the corpus level under the tested conditions (single corpus, single classifier, binary-gender lexicon, quick-run scale). Whether this restructuring propagates to model behavior and downstream fairness outcomes remains an open question requiring full-scale training experiments.

\textbf{Directions for future work include:} full-scale Pythia-1B training at 100B tokens with gpt-neox framework to test the H-M2 graded correlation; execution of H-M3 (fairness benchmark comparison on matched-capability models); extension of the corpus audit methodology to other corpora (RefinedWeb, FineWeb, RedPajama); a continuous 20-level fastText percentile sweep to test the functional form of the log-odds relationship; and expansion of the demographic lexicon beyond binary-gender, English-only, U.S.-centric categories.

